\section{Carrera e investigación: crear novel constraints en sistemas reales}

Un estudiante le pregunta a Rauch si se arrepiente de haber entrado a la industria demasiado pronto en lugar de seguir la ruta de investigación tradicional. Su respuesta no es «la industria es mejor que la academia», sino que él mismo participó en la empresa en trabajos de tipo investigación, como el scheduler, el placement, el load balancer y el metadata store. Un producto real puede ofrecer gran escala, restricciones estrictas y retroalimentación inmediata; pero también reconoce que hacerlo todo uno mismo en exceso (DIY) puede acabar con una startup.

\begin{knowledgebox}{Tres condiciones necesarias para la Applied research}
\begin{enumerate}
  \item El problema contiene una constraint explícita que los métodos existentes no pueden satisfacer;
  \item el equipo puede construir experimentos, mediciones y retroalimentación de fallas, en lugar de solo escribir parches de producción;
  \item los resultados de la investigación pueden volver a las métricas del producto y, al mismo tiempo, conservar conocimiento reutilizable y límites del sistema.
\end{enumerate}
\end{knowledgebox}

\begin{warningbox}{«Escribirlo uno mismo» no es la ruta de aprendizaje predeterminada}
Reescribir React, Kubernetes o un scheduler puede generar una comprensión profunda, pero también puede desperdiciar el runway de la empresa. Conviene elegir la capa que esté directamente relacionada con la diferenciación central y cuyas restricciones no puedan satisfacer las soluciones existentes; para el resto, hay que dar prioridad a la reutilización y aprender mediante la lectura del código fuente, la experimentación y los simulacros de fallas.
\end{warningbox}

\teachervoice{La verdadera tensión de la clase está en que dos frases son ciertas a la vez: too much DIY can kill you; unreasonable standards can create novel research. Un criterio maduro no consiste en «comprarlo todo» ni en «construirlo todo», sino en saber qué restricción merece que la organización invierta su vida en resolverla.}

\subsection{Resumen de la sección}

Los sistemas industriales pueden convertirse en un entorno de investigación, siempre que las restricciones del problema, la retroalimentación experimental y el valor para el producto estén claros. Build vs. Buy se aplica también al aprendizaje personal: profundizar en las capas clave y reutilizar las capas que no lo son.
